\fancyhead[L]{\small BRIEF AN BRANDES.}
\SourceLeaf{515}\pdfbookmark[0]{Brief von Gauss an Brandes}{brief_brandes}
\expandafter\def\csname HistoricalNoteText515\endcsname{[Über die achromatischen Doppelobjektive, besonders im Hinblick auf die vollständige Aufhebung der Farbzerstreuung.]}
\begin{center}
\hrule\vspace{2mm}
\textsc{Gehler's} Physikalisches Wörterbuch. 1831. Artikel:\par
Linsenglas, Berechnungen über achromatische und aplanatische Linsengläser aus zwei Glaslinsen.
\vspace{2mm}\hrule\vspace{5mm}
{\large\itshape Brief von Gauss an Brandes.\par}
\end{center}
Ihr Brief hat mich veranlasst, eine freie Stunde auf den Umstand zu verwenden, der am Ende jenes Aufsatzes\HistoricalNote{515} kurz erwähnt wird. Der eigentliche Sinn der dortigen Bemerkung scheint nicht von allen ganz richtig verstanden worden zu sein, aber auch meine eigene Angabe bedarf einer kleinen Änderung. Ich finde nämlich jetzt durch eine \textit{eingehendere} Untersuchung, dass die Unschärfe, die im Ausdruck für die Längsabweichung von der vierten Potenz des Abstands der einfallenden Strahlen von der Achse abhängt, den kleinstmöglichen Gesamteinfluss hat, wenn man das Objektiv so konstruiert, dass die unendlich nahe an der Achse einfallenden Strahlen und diejenigen, die in einem Abstand $=R.\BareRoot{\tfrac65}$ einfallen würden (wobei $R=$ der Radius des Objektivs ist), sich in \textit{einem} Punkt $A$ vereinigen. Das Okular ist dann so eingestellt, dass man den Punkt der Achse deutlich sieht, in dem sich die Strahlen, die in den Abständen $=(\tfrac35-\tfrac{\BareRoot{6}}{10})R$ und $=(\tfrac35+\tfrac{\BareRoot{6}}{10})R$ von der Achse eingefallen sind, alle vereinigen. Denken Sie sich durch diesen Punkt eine Ebene senkrecht zur Achse. Dann ist das Bild umso unschärfer, je größer der Kreis um $A$ ist, den die von einem Punkt des Gegenstands auf das Objektivglas gefallenen Strahlen ausfüllen. Dabei muss die Intensität der Strahlen an jeder Stelle des Kreises mit berücksichtigt werden. Hier liegt eine gewisse Willkür vor. Ich halte es für am zweckmäßigsten, nach denselben Prinzipien vorzugehen, die der Methode der kleinsten Quadrate zugrunde liegen. Ist $\mathrm{d}s$ ein Element dieses Kreises, $\rho$ der Abstand des Elements von $A$ und $i$ die Intensität der Strahlen an dieser Stelle, so nehme ich an, dass $\int i\rho\rho\,\mathrm{d}s$ als Maß der Gesamtunschär\SourceLeaf{516}fe gelten kann, und mache es zu einem Minimum. Dabei finde ich folgende Ergebnisse. 1. Man konstruiere das Objektiv so, dass das vom Quadrat des Abstands von der Achse abhängige Glied der Längsabweichung $=0$ wird, und stelle das Okular so. dass $A$ dorthin fällt, wo die der Achse unendlich nahen Strahlen diese schneiden. Der Wert dieses Integrals sei dann $=E$. 2. Stellt man dagegen \textit{bei derselben Anordnung} das Okular so ein, dass das Integral so klein wird, wie es bei dieser Anordnung werden kann (wobei $A$ der Vereinigungspunkt der in einem Abstand $=R\BareRoot[4]{\tfrac12}$ einfallenden Strahlen sein wird), so ist das Integral $=\tfrac14E$. 3. Bei der oben beschriebenen Anordnung und der vorteilhaftesten Einstellung des Okulars ist das Integral dagegen $=\tfrac{1}{100}E$, sein absolutes Minimum. Das obige Ergebnis, dass nämlich ein bloß \textit{gedachtes} Bild (von Strahlen aus einem größeren Abstand von der Achse als dem Radius des Objektivs) mit dem Vereinigungspunkt der unendlich achsennahen Strahlen zusammenfallen soll, ist zunächst sehr überraschend und scheint paradox. Bei näherer Betrachtung erkennt man den eigentlichen Grund aber leicht. Jenes erste \textit{sogenannte} Hauptbild (von Strahlen sehr nahe an der Achse) ist dabei wegen seiner geringen Intensität gewissermaßen das Unwichtigste. Viel wichtiger ist, dass die Strahlen von den näher am Rand des Objektivs gelegenen Ringen \textit{untereinander} besser zusammengehalten werden. Das wird bei jener Anordnung am besten erreicht. Ich bedaure, dass die Grenzen eines Briefes jetzt keine ausführlichere Darstellung erlauben. Gegen die genaue Rechnung lässt sich nichts einwenden, und bei einer vagen Überlegung übersieht man leicht einen wesentlichen Umstand. Für den Kenner werden diese Hinweise genügen.

Allgemein finde ich, dass \textit{bei der vorteilhaftesten Einstellung des Okulars} jenes Integral stets $=\tfrac14E(1-\tfrac85\mu\mu+\tfrac23\mu^4)$ wird, wenn das Objektiv so konstruiert ist, dass sich die Strahlen aus dem Abstand $\mu R$ von der Achse mit dem oben so genannten Hauptbild in einem Punkt vereinigen. Das ist ein Minimum für $\mu=\BareRoot{\tfrac65}$ und dann $=\tfrac{1}{100}E$. Für $\mu=1$ wäre es nur $=\tfrac{1}{60}E$ und für $\mu=$ unendlich klein, $=\tfrac14E$. Demnach hat \textsc{Bohnenberger} Unrecht, und auch ich hatte damals Unrecht, insofern als ich noch nicht weit genug von \textsc{Bohnenberger} abgewichen war. Ich hatte damals nur die \textit{gesamte Größe} des unscharfen Bildes berücksichtigt, ohne auf die ungleiche Intensität der einzelnen Teile zu achten.

